\section{Introduction}\label{sec:introduction}

For a smooth integral projective variety $X$, the pluricanonical
systems organize the birational information carried by its canonical
divisor. If some $H^0(X,\OO_X(mK_X))$ is nonzero, the maximum of the
dimensions of the resulting rational images is the Kodaira dimension
$\kappa(X)$. In sufficiently large divisible degrees these maps determine
the Iitaka fibration \cite{Iitaka1971}. The question here is whether divisibility and
largeness can be prescribed using only $\dim X$.

To specify the required conclusion, let $K(K_X)\subset k(X)$ be the field
generated by ratios of nonzero sections in the same positive
pluricanonical degree, over all such degrees. This is the function field
of the Iitaka base. A complete system \emph{defines the Iitaka fibration}
if it is nonempty and its section ratios generate $K(K_X)$. Requiring
only an image of dimension $\kappa(X)$ would be weaker: it would allow
a proper finite subfield of the Iitaka field.

\begin{theorem}[Uniform pluricanonical degree]\label{thm:main}
For each integer $d\geq1$ there is an integer $m(d)>0$ such that, over
every algebraically closed field of characteristic zero, the complete
system $|m(d)K_X|$ defines the Iitaka fibration of every smooth integral
projective $d$-fold $X$ with $\kappa(X)\geq0$.
\end{theorem}

Using the log-abundance input from \cite{LA}, this gives an affirmative
resolution of the effective Iitaka fibration conjecture
\cite[Conjecture~1.1]{BirkarZhang2016}. Every positive multiple of
$m(d)$ works as well, by the section-ratio argument in
Section~\ref{sec:completion}.
Thus the same degree works for all nonnegative Kodaira dimensions. In
the general-type case its map is birational. In Kodaira dimension zero
the conclusion is nonemptiness in that degree, with image a point.
The theorem is existential; no small or practical numerical value of
$m(d)$ is asserted.

The proof establishes a second uniform statement at the same time.
A pair in the next theorem is normal and integral; its boundary is
effective and its log canonical divisor is rational Cartier.

\begin{theorem}[Uniform log canonical index]\label{thm:lc-index}
Let $d\geq0$, and let $\Phi\subset[0,1]\cap\Q$ satisfy the descending
chain condition. There is an integer $a(d,\Phi)>0$ such that every
projective log canonical pair $(X,B)$ of dimension $d$, with boundary
coefficients in $\Phi$ and $K_X+B\sim_\Q0$, satisfies
\[
                         a(d,\Phi)(K_X+B)\sim0.
\]
Here the multiple is an integral principal divisor, not merely a
numerically trivial or rationally linearly trivial divisor.
\end{theorem}

\subsection{Context and preceding work}

Effective pluricanonical maps are a quantitative part of birational
classification. The general-type case asks for a uniform birational
degree and was established by Hacon--McKernan, Takayama, and Tsuji
\cite{HaconMcKernan2006,Takayama2006,Tsuji2007}.
Hacon--McKernan--Xu's theorem for big log canonical adjoints permits
boundary coefficients in a fixed DCC set \cite[Theorem~1.3]{HMX2014}.
The lower-Kodaira-dimensional problem has an additional feature:
the general fibre contributes torsion and monodromy data which do not
appear in the big case.

Fujino--Mori obtained a uniform pluricanonical degree for threefolds
of Kodaira dimension one using their canonical bundle formula
\cite[Corollary~6.2]{FujinoMori2000}. Their denominator argument
identifies finite characters on an extended Hodge line and bounds
their orders using the middle Betti number of a canonical cover
\cite[Theorem~3.1 and Sections~3.6--3.8]{FujinoMori2000}.
Viehweg--Zhang obtained effective results for Iitaka bases of dimension
two, with constants controlled by invariants of the general fibre
\cite[Theorem~0.2]{ViehwegZhang2009}. Birkar--Zhang proved a general
effective Iitaka theorem with dependence on the dimension, the least
nonvanishing canonical degree of the general fibre, and a middle Betti
number of a smooth model of its canonical cover
\cite[Theorem~1.2]{BirkarZhang2016}. Their effective birationality theorem
for polarized adjoints is the final effective input used here
\cite[Theorem~1.3]{BirkarZhang2016}. The present argument supplies the
required denominator control without assuming a uniform bound for
those middle Betti numbers.

The index problem has its own inductive structure. Xu relates
zero-boundary klt index bounds to bounds with boundary
\cite[Theorem~1.8]{Xu2019}. Birkar proves the index theorem for rationally
connected klt Calabi--Yau pairs in every fixed dimension
\cite[Corollary~1.8]{Birkar2025}; his boundedness theorem for generalized
Calabi--Yau pairs also controls torsion on rationally connected bases
\cite[Theorem~1.7]{Birkar2025}. We combine these results with an arithmetic
Stein-degree theorem \cite[Main Theorem]{SD} to pass from klt pairs to
normal lc pairs without an index theorem for nonnormal pairs.

\subsection{Inputs and contributions}

One major input must be distinguished from standard minimal-model
technology. We use the good-log-minimal-model theorem of the companion
\emph{Log abundance in characteristic zero}
\cite[Theorem~11.1]{LA}, stated precisely in
Theorem~\ref{thm:good-models}. It includes nonvanishing and is used in
all dimensions. It is not being inferred from finite generation for
big adjoints. The other companion inputs are the chain and tracking
lemmas in \cite{U4}, and the arithmetic Stein-degree bound in \cite{SD}.
The arguments of \cite{R4} provide the preceding fourfold treatment of
relative denominators and rationally connected torsion.

The proof extends the canonical-index method of \cite{U4} to arbitrary
dimension. There are two substantial extensions. First, we prove a
relative denominator bound by decomposing the geometric generic fibre,
degenerating its individual factors, and controlling their volume-form
characters on normal components of a dlt special fibre. This separates
the potentially unbounded ramification of semistable reduction from
the bounded characters that enter the moduli divisor. The argument
uses the absence of horizontal boundary and only lower-dimensional
normal lc index bounds. It is therefore a reusable relative statement,
not merely the final step of the Iitaka theorem.

Second, the small-volume argument in \cite{U4} is extended to Iitaka
bases of arbitrary dimension. Weak positivity supplies a fibre-to-total
volume comparison, while flattening ensures that codimension-two
losses on the base create no divisorial poles on the original smooth
model. This gives scalar bounds for section orders and Seshadri
constants. We then give the dimension-dependent calculations in the
Frobenius diagonal argument and in the compatible-chain argument.
Those constructions are adaptations of \cite{U4}; the new claim is
their use in the all-dimensional induction, not a new origin for the
constructions themselves.

\subsection{The induction}

We prove Theorems~\ref{thm:main} and~\ref{thm:lc-index} simultaneously.
Fix $n$ and suppose both are known in smaller dimensions. The relative
step starts with a contraction $f:X\to Z$ from a klt $n$-fold, with
$\dim Z>0$ and $K_X$ rationally linearly pulled back from $Z$.
The canonical bundle formula writes
\[
                    K_X\sim_{\Q}f^*(K_Z+B_Z+M_Z),
\]
where $B_Z$ is the discriminant divisor and $M_Z$ is the moduli
divisor. Section~\ref{sec:denominators} bounds a Cartier multiple of
the moduli divisor on a birational base model using only
lower-dimensional normal lc indices. When $K_X\sim_{\Q}0$ and the
base has a rationally connected smooth resolution, this denominator
bound combines with torsion control on the base to bound the canonical
index of $X$.

Section~\ref{sec:reductions} then reduces failure of the zero-boundary
klt index bound to terminal varieties $V$ with canonical indices
$r\to\infty$ and countable birational groups. Their canonical cyclic
covers $\pi:Y\to V$ have degree $r$ and deck group $G$. Every
$G$-equivariant rational fibration of $Y$ with positive-dimensional
base and fibre has general-type smooth geometric generic fibre.
We call such fibrations intermediate. The corresponding assertion holds for
intermediate rational fibrations of $V$.

Choose an ample rational Cartier divisor $L$ on $V$ with
$1\le L^n\le2$. For an ample rational Cartier divisor $P$, write
$\varepsilon(P)$ for its very general Seshadri constant: the infimum
of $P\cdot C/\operatorname{mult}_xC$ over integral curves through a
very general point $x$. Section~\ref{sec:scalar} uses lower-dimensional
effective Iitaka fibrations to bound section orders on $V$ and bound
$\varepsilon(L)$ below uniformly. The pullback polarization has
$(\pi^*L)^n=rL^n$.

For $t\ge2$, let $Z_t=Y^t/G_{\mathrm{diag}}$, where $G$ acts by
the same deck transformation in every slot. Let
$\theta_t:Z_t\to V^t$ be the induced finite map and set
$P_t=\theta_t^*(\sum_{i=1}^t\operatorname{pr}_i^*L)$.
Together, the scalar estimates and the diagonal argument of
Section~\ref{sec:diagonal} give
\[
       \varepsilon(\pi^*L)\ge c r^{1/n},
       \qquad \varepsilon(P_t)\le Ct^2,
\]
with $c,C>0$ depending only on $n$. The second bound supplies curves
of bounded $P_t$-degree divided by marked multiplicity.

Section~\ref{sec:transport} arranges these curves into chains compatible
with coordinate projections. The resulting positive-dimensional
leaves project generically finitely onto their images in each
coordinate, and comparison of their projection degrees gives a
degree-one coordinate-forgetting map.
Rational recovery of its missing coordinate turns local flows into
uncountably many birational self-maps of $V$, a contradiction. This
bounds the zero-boundary klt index. The normal lc reduction in
Section~\ref{sec:index} then proves Theorem~\ref{thm:lc-index} in
dimension $n$. Finally, Section~\ref{sec:completion} reuses the
denominator bound with polarized effective birationality to prove
Theorem~\ref{thm:main} in that dimension. Figure~\ref{fig:induction}
records the dependencies after the induction statements have been
formally introduced.
